%%%PAGE 100 rot=0 legibility=fair summary=机械振动与机械波知识提纲(思维导图式)：简谐运动两种模型、公式图像、受迫振动与共振、规律与特征、单摆周期公式、机械波基本概念、波的图像与多解问题
%%%BLOCK id=100-01 ch=08 sec=简谐运动·两种模型 kind=summary alt=none cont=none y=0.09-0.28
% 原稿：左侧一个大括号并列“简谐运动”与“简谐运动的两种模型”；“两种模型”右侧再用大括号括出下面两项
\begin{itemize}
\item 简谐运动
\item 简谐运动的两种模型
  \begin{itemize}
  \item 弹簧振子（水平）
  \item 单摆
  \end{itemize}
\end{itemize}
自由振动\quad $f=f_{\text{固有}}$\quad $T=T_{\text{固有}}$
% “自由振动”一行位于右大括号内、“弹簧振子（水平）”与“单摆”之间偏右处
%%%END
%%%BLOCK id=100-02 ch=08 sec=简谐运动·公式与图像 kind=summary alt=none cont=none y=0.24-0.34
% 以下各项（至“单摆及周期公式”）在原稿中被一个大括号归并在左侧“机械振动”之下
\textbf{机械振动。}
\begin{itemize}
\item 简谐运动的公式和图像。\quad $F_{\text{回}}=-kx$（回复力\ 效果力）指向平衡位置。
  \begin{itemize}
  \item $\uparrow$正余弦曲线。\quad $T=2\pi\sqrt{\dfrac{L}{g}}$\quad $T=2\pi\sqrt{\dfrac{m}{k}}$
  \end{itemize}
\end{itemize}
% “正余弦曲线”上方有向上箭头指向“图像”
%%%END
%%%BLOCK id=100-03 ch=08 sec=受迫振动与共振 kind=summary alt=none cont=none y=0.27-0.43
\begin{itemize}
\item 受迫振动与共振\quad 机械能不守恒
  \begin{itemize}
  \item 系统在驱动力下的振动。
  \item 受迫振动频率等于驱动力频率，当驱动力频率等于系统固有频率时，发生共振。
  \end{itemize}
\end{itemize}
% 原稿：“机械能不守恒”写在“受迫振动与共振”上方偏右处；其左侧另有两行小字如下
机械工作底座振动\\
共振筛、\unclear{声呐}

（共振曲线：纵轴$A$（受迫振动振幅），横轴$f_{\text{驱}}$，峰值对应$f_{\text{固}}$。）
%FIG p100_a 0.775 0.337 0.955 0.406
\notefig[0.28]{figures/p100_a.png}
%%%END
%%%BLOCK id=100-04 ch=08 sec=简谐运动·规律与特征 kind=summary alt=none cont=none y=0.38-0.50
\begin{itemize}
\item 简谐运动的规律。
  \begin{itemize}
  \item 受力特征\quad $F=-kx$
  \item 运动特征：$a$ $F$ $x$ 同\unclear{增}减，$v$ 与其相反
  \item 能量特征
    \begin{itemize}
    \item 振幅越大能量越大，
    \item 动能势能相互转化，机械能守恒
    \end{itemize}
  \item 对称性特征
    \begin{itemize}
    \item 关于$O$对称，$aFx$大小相等的\unclear{相应}，$V$大小相等，方向可同可反，动能势能相等
    \end{itemize}
  \item 周期性特征
    \begin{itemize}
    \item 标量为$\dfrac{1}{2}$\unclear{$T$}，矢量为$T$
    \end{itemize}
  \end{itemize}
\item 简谐运动的理解及应用
  \begin{itemize}
  \item 图像的信息利用
  \end{itemize}
\end{itemize}
%%%END
%%%BLOCK id=100-05 ch=08 sec=单摆·周期公式 kind=formula alt=none cont=none y=0.50-0.58
\begin{itemize}
\item 单摆及周期公式。\quad $F_{\text{回}}$\ $F_{\text{向}}$\ 大小方向时刻变。\quad $T=2\pi\sqrt{\dfrac{L_{\text{等效摆长}}}{g_{\text{等效重力加速度}}}}$
\end{itemize}
\[ F_{\text{回}}=-mg\sin\theta=-mg\frac{x}{L}=-kx\qquad F_{\text{向}}=T-mg\cos\theta \]
%%%END
%%%BLOCK id=100-06 ch=08 sec=机械波·基本概念 kind=summary alt=none cont=none y=0.57-0.77
% 以下各项在原稿中被一个大括号归并在左侧“机械波”之下；左侧三行标注如下
机械波\\
传播形式和能量\\
不传播质点。
\begin{itemize}
\item 机械波。横波和纵波\ $\leftarrow$\ 振动方向与传播方向平行，有密部和疏部。
\item 波源$+$介质\quad $\hookrightarrow$\ 振动方向与传播方向垂直
\item 横波的图像。波长、波速和频率的关系。\quad $v=\lambda f$\quad 同介质，$v$ 相同。
\item 介质中各点 $T$、$A$、$f$ 都与波源相同。一个周期路程 $4A$，位移为0。\quad $\lambda=vT$
\item 波的干涉和衍射，多普勒效应\ 波源与观察者间有相对运动，波源与波\unclear{频}率不变，观察者接收到的频率\unclear{增大}
  \begin{itemize}
  \item $\hookrightarrow$\ 两波频率必须相等。
  \item $\hookrightarrow$\ 波长与障碍物或孔的尺寸差不多或大。
  \end{itemize}
\end{itemize}
%%%END
%%%BLOCK id=100-07 ch=08 sec=波的图像与多解问题 kind=summary alt=none cont=none y=0.77-1.0
\begin{itemize}
\item 机械波的传播与波的图像。
  \begin{itemize}
  \item $y$-$x$ 图
  \end{itemize}
\item 振动图像与波的图像综合应用
  \begin{itemize}
  \item 用时间联系
  \end{itemize}
\item 波的多解问题。
  \begin{itemize}
  \item 双向性
    \begin{itemize}
    \item 传播方向双向性
    \item 振动方向双向性
    \end{itemize}
  \item 周期性
    \begin{itemize}
    \item $\Delta t$ 与 $T$ 关系不确定
    \item $\Delta x$ 与波长关系不确定。
    \end{itemize}
  \item 隐含性\quad 是只给出波形的一部分或只给出几个特殊点。
  \end{itemize}
\item 波的干涉、衍射和多普勒效应
\end{itemize}
%%%END
%%%PAGEEND 100
